% Typeset transcription proof: source journal page 293 (PDF page 2).
\documentclass[11pt]{article}
\usepackage[T1]{fontenc}
\usepackage[utf8]{inputenc}
\usepackage[margin=1.05in]{geometry}
\usepackage{amsmath,amssymb}
\usepackage{mathptmx}
\setlength{\parindent}{1.5em}
\setlength{\parskip}{0pt}
\begin{document}
\thispagestyle{empty}

\begin{center}
  \textit{Multenions and Differential Invariants.}\hfill 293
\end{center}

Euclidian space of $n$ dimensions, $i_1$, $i_2$, \ldots{} $i_n$. These are
such that all the ordinary laws of algebra except the commutative law of
multiplication apply to the general integral expression $q$, called a
multenion, where

\begin{equation}
q = \sum (x_0+x_1i_1+x_2i_2+\cdots)(x'_0+x'_1i'_1+x'_2i'_2+\cdots)\cdots,
\tag{1}
\end{equation}

where every $x$ is a scalar. Since any $x$ may be negative we understand this
statement to include the ordinary algebraic uses of the signs $+$ and $-$; but
we do not understand that any assertion has been made about the use of the sign
$\div$ [$q\div r$ can later be understood to be an alternative for $qr^{-1}$].

\noindent\textsc{Law II.} Scalars are commutative not only with each other but
with each of the $n$ primitive unit vectors. [Thus scalars are commutative with
multenions.]

\begin{center}
  \textsc{Law III.} [$i_1^2$ means $i_1i_1$, etc.]
\end{center}
\begin{equation*}
  -1=i_1^2=i_2^2=\cdots=i_n^2.
\end{equation*}

\noindent\textsc{Law IV.} In a product, called a primitive vectorium,
$i_ai_b\ldots$, of primitive unit vectors, if two adjacent \emph{different}
vectors be interchanged, the product merely changes sign.

\indent Such are the assumptions. The rest is development, aided of course by
suitable terminology and symbolisation.

\noindent\textsc{Def. 1.} A primitive unit means any one of the following:
(1) the scalar $1$, (2) any primitive unit vector $i_a$, (3) any other
primitive vectorium $i_ai_b\ldots$. [Strictly (3) includes (2) and in a sense
(1).]

\noindent\textsc{Def. 2.} An expression of the form
\begin{equation*}
  p=\sum_{c=1}^{n}x_ci_c
\end{equation*}
where each $x$ is a scalar, is called a vector, or a $V_1$ or a ${}_nV_1$.

\noindent\textsc{Def. 3.} Any multenion is obviously capable of expression in
the form $\sum xi_ai_b\ldots i_k$ where $i_a$, $i_b$, \ldots{} $i_k$ are $m$
different primitive vectors and $m$ ranges from $0$ to $n$, $m$ being given
$V_mq$ means the sum of all parts of the multenion for that value of $m$. Thus

\begin{equation*}
  q=V_0q+V_1q+V_2q+\cdots+V_nq.
\end{equation*}

$V_aq$ is said to be the part of $q$ of homogeneity $a$ and is called a $V_a$
or a ${}_nV_a$ or a homogeneous multenion or a hyper vector. [These alternatives
are intended to invite suggestions for a settlement of the name. I should have
preferred ``an $a$-vector'' but this clashes with relativity meanings of
4-vector and 6-vector. In multenions with $n=4$, there are two kinds of
4-vector ${}_4V_1$ and ${}_4V_3$.]

$V_0q$ will generally in words be called the scalar part of $q$ but in the
present connection we must not use $Sq$ as an alternative for $V_0q$. In
applications we very decidedly need to be able, in one system of multenions
($n=4$) to make

\newpage
\thispagestyle{empty}
\begin{center}
  294\hfill Prof.\ A.\ McAulay.
\end{center}

no change in the quaternion notation but yet to use both the quaternion
notation and the full multenion notation as applying to one symbol $q$. Thus
\begin{equation*}
  Sq=V_0q+V_4q, \qquad Vq=V_2q.
\end{equation*}
Endless confusion would result from allowing $V_0$ to be represented by $S$.

\noindent\textsc{Def. 4.} If $\alpha_1$, $\alpha_2$, \ldots{} $\alpha_a$ are
any $a$-vectors then $V_a\alpha_1\alpha_2\ldots\alpha_a$ is called a
vectorium. [Hence we have already called such an expression as $i_1i_2i_3$ a
primitive vectorium.]

\noindent\textsc{Def. 5.} A linear multenion function of a multenion is called
a multenion linity. Similarly for a vector linity or a ${}_nV_a$ linity, etc.
An extended vector linity is a special kind of multenion linity, $\Phi$, defined
from a given vector linity, $\phi$. Its fundamental property is this: If
$\alpha_1$, $\alpha_2$, \ldots{} $\alpha_a$ are any $a$-vectors whatever, $a$
itself ranging from $0$ to $n$, then
\begin{equation}
  \Phi V_a\alpha_1\alpha_2\ldots\alpha_a
    = V_a\phi\alpha_1\phi\alpha_2\ldots\phi\alpha_a.
\tag{2}
\end{equation}

[It is not obvious that this remark covers a possibility in general. The
definition, which will be modified below, is given here as a hint of the
importance of vectoriums and extended linities in our subject.]

\begin{center}
  \textit{Usual Notations.}
\end{center}

\noindent Scalars ($V_0$): $a$, $b$, $c$, $g$, $h$, $k$, $l$, $m$, $n$,
$t$ ($=$ time), $x$, $y$, $z$, $\mathfrak{J}$ ($=\sqrt{-1}$), $\theta$
($=$ arc), $\pi$.

\noindent General multenions: $p$, $q$, $r$, $s$.

\noindent Hyper vectors (homog. mult.) and vectoriums:
\[
\begin{array}{|l|c|c|c|c|c|}\hline
\text{Homogeneities} & \substack{V_a\\a} & \substack{V_b\\b} & \substack{V_c\\c} & \substack{V_n\\n} & \substack{V_2\\2}\\ \hline
\text{Hyper vectors} & u & v & w & x\upsilon & \omega \\[2pt]
\text{Vectoriums} & \breve{\alpha} & \breve{\epsilon} & \breve{\eta} & x\breve{\upsilon} & \breve{\omega} \\ \hline
\end{array}
\]

\noindent (These conventions should be noted and remembered. They immensely
simplify the reading of formulae by their property of minimising number of
suffixes, etc.)

\noindent Vectors ($V_1$): $\alpha$, $\beta$, $\gamma$, $\delta$, $\epsilon$,
$\zeta$, $\xi$, $\eta$, $\lambda$, $\mu$, $\nu$, $\rho$, $\sigma$, $\tau$.

\noindent Linear functions: $\xi$, $\eta$, $\lambda$, $\mu$, $\nu$, $\varpi$,
$v$, $\phi$, $\chi$, $\psi$.

\noindent Vectoriums: $\breve{\alpha}$, $\breve{\epsilon}$, $\breve{\eta}$,
$\breve{l}$, $\breve{\sigma}$, $\breve{\beta}$, $\breve{\upsilon}$, $\breve{\omega}$.

\noindent Differential symbols: $\delta$, $\Delta$, $\mathrm{H}$, $\mathrm{L}$,
$d$, $d_\rho$, $d_{\breve{\rho}b}$, $d_{s\alpha}$, $d\alpha$, $d\breve{b}$.

\noindent Selective linities, etc.: $\mathrm{K}$, $\mathrm{P}$, $\mathrm{Q}$,
$\mathrm{R}$, $\mathrm{S}$, $\mathrm{T}$, $\mathrm{U}$, $\mathrm{V}$.

\noindent\textsc{Def. 6.} Complexes of vectors, hypervectors, multenions,
etc., are sufficiently illustrated by multenions. If $q_1$, $q_2$, \ldots,
$q_a$ are $a$ given multenions, then all multenions of the form
$\sum_{c=1}^{a}x_cq_c$ are said to form a complex. If $q_1$, $q_2$, \ldots,
$q_a$

\newpage\thispagestyle{empty}
\begin{center}\textit{Multenions and Differential Invariants.}\hfill 295\end{center}
can be made linearly dependent on $b$ and no fewer, then $b$ is called the complexity of the complex. Thus the complexity of the complex of all vectors is $n$, of all multenions it is $2^n$, of all multenions of homogeneity $a$ it is ${}_nC_a$.

\noindent\textsc{Def. 7.} A unit vector or a unit vectorium is one, $\breve a$, such that $\breve a^2=\pm1$. [The square of every vectorium is a scalar, but $u^2$ is a scalar in general only in the four cases when $a=0,1,n-1$ or $n$.] Thus all the primitive vectoriums are unit vectoriums.

\noindent\textsc{Def. 8.} Two vectors, $\alpha$, $\beta$, are said to be normal to one another when $V_0\alpha\beta=0$. Two vector complexes are said to be completely normal to one another when every vector of the one complex is normal to every vector of the other.

\indent Every one of the symbols $V_a$ is a multenion linity. They are all commutative with one another. They are all idempotent, that is $V_a^2=V_a$, and they all satisfy the equation $V_aV_b=0$ when $a\ne b$.

\indent There is another system of multenion linities of great utility and simplicity. They are $2^n$ in number. They are all commutative with each other and with each of the $n+1$ linities $V_a$. Let $P_a$ be that definite linity of a general multenion, $q$, which simply reverses the sign of every $i_a$ which occurs in the expression (1) for $q$. From the well-known properties of linities a product of the P's such as $P_1P_2P_3$ simply reverses the sign of every $i_1$, $i_2$, and $i_3$. These linities are (1) unipotent, that is, if
\begin{equation}\phi=P_aP_b\cdots P_k,\tag{3}\end{equation}
then $\phi^2=1$; (2) they are proscriptive, by which is meant that
\begin{equation}\phi(qr)=\phi q\,\phi r,\tag{4}\end{equation}
where $q$ and $r$ are any two multenions. It is convenient to use $P$ for the product $P_1P_2\ldots P_n$ of all the linities $P_a$, so that $P$ reverses the sign of every primitive vector.

\indent $K$ is defined, consistently with its quaternion use, in the following manner. $Kq$ is a linity of a general multenion, $q$, which reverses the sequence of the primitive unit vectors in every primitive vectorium $\brevel$, of $q$, and \emph{also} reverses the sign of every primitive unit vector. Thus $K(i_1i_2i_3)=(-i_3)(-i_2)(-i_1)=i_1i_2i_3$. $K$ is (1) unipotent, that is $K^2=1$; (2) it is commutative with every $V$ and every $P$; (3) it is retrospective, that is $K(qr)=Kr\,Kq$. [To prove (3) first prove it when $q$ and $r$ are two primitive vectoriums and then generalise. Such a method of proof is very often applicable and may be referred to as proof by reduction to primitive units.]
\begin{equation}K(qr)=Kr\,Kq.\tag{5}\end{equation}

\newpage\thispagestyle{empty}\begin{center}296\hfill Prof.\ A.\ McAulay.\end{center}
but if $\phi$ be given by (4), $K\phi(=\phi K)$ is also retrospective. More generally, if $\phi$, $\psi$, $\chi$, \ldots{} are any number of multenion linities, every one of which is either proscriptive or retrospective, then also is the product, $\phi\psi\chi\ldots$, proscriptive or retrospective; it is proscriptive when the number of $\phi$, $\psi$, $\chi$, \ldots{} which are retrospective is even, and it is retrospective when that number is odd.

\indent The special retrospective $PK(=KP)$ is of sufficient importance to have a symbol appropriated to it. We shall call it $Q$. Thus the product of any two of the three $P$, $Q$, $K$ is equal to the third.

\indent A caution to users of present methods in relativity and hyperbolic geometry seems desirable. It relates to the discussion of real questions by aid of the imagery
\begin{equation}\mathfrak J=\sqrt{-1}.\tag{6}\end{equation}
\indent Let us illustrate by anticipating our mode of treatment of real questions in relativity. We begin with the multenion system based on the primitive generators $i_1$, $i_2$, $i_3$, $i_4$. Then we put
\begin{equation}i_4=\mathfrak Jl,\tag{7}\end{equation}
and drop out of most of our subsequent work all reference to the original $i_4$ and $\mathfrak J$, using $l$ instead with its property $l^2=1$, as opposed to $i_4^2=-1$. Questions in relativity are then treated by aid of the multiplicative combinations of such expressions as
\begin{equation}xi_1+yi_2+zi_3+tl.\tag{8}\end{equation}
whereby, so long as all the scalars thereby introduced are real, we shall be dealing with real questions in relativity. The general multenion thus constructed contains 16 and not 32 independent real scalars. $l$ is for practical relativity purposes real. We shall call such a multenion system semi-real.

\indent Now return to the meaning above assigned to $K$.
\begin{align*}K(i_1i_2i_3i_4)&=i_1i_2i_3i_4,\\ \text{and also}\quad K(i_1i_2i_3l)&=i_1i_2i_3l,\\ \text{but whereas}\quad K'(i_1i_2i_3i_4)&=(i_1i_2i_3i_4)^{-1}.\tag{9}\end{align*}
it is not true that the corresponding equation holds when we replace $i_4$ by $l$. I have therefore refrained from making the definition of $K$ depend on such equations as (9). It is, nevertheless, a perfectly easy theorem to prove that
\begin{equation}Kl=l^{-1}.\tag{10}\end{equation}
when $l$ is any product of our \emph{original} generators $i_1$, $i_2$, \ldots{} $i_n$. This ceases to be true when $l$ takes the place of $i_4$, simply because $\mathfrak J^{-1}$ is $-\mathfrak J$ and not $+\mathfrak J$. It may be noted, however, that there is always a reproscriptive $K\phi=R$ to be found, which is such that for every semi-real vectorium $\brevel$, $R\brevel=\brevel^{-1}$.

\newpage\thispagestyle{empty}
\begin{center}\textit{Multenions and Differential Invariants.}\hfill 297\end{center}
Thus, suppose we replace $i_{a+1}$, \ldots{} $i_n$ by $l_{a+1}$, \ldots{} $l_n$, where
\begin{equation}i_c=\mathfrak J l_c\quad\text{and therefore}\quad l_c^2=1,\qquad c=a+1,a+2,\ldots,n.\tag{11}\end{equation}
Then
\begin{equation}R=KP_{a+1}P_{a+2}\cdots P_n=QP_1P_2\cdots P_a,\tag{12}\end{equation}
is easily shown to have the required characteristic. We now obviously have
\begin{align}P&=P_1P_2\cdots P_n,\tag{13}\\ Pu&=(-1)^a u,\qquad Qu=(-1)^{\frac12a(a-1)}u,\qquad Ku=(-1)^{\frac12a(a+1)}u.\tag{14}\end{align}
Further if we put
\begin{equation}q_c=V_cq,\qquad c=0,1,\ldots,n,\tag{15}\end{equation}
\begin{align}Pq&=q_0-q_1+q_2-q_3+\cdots,\notag\\P_{(0)}q&=\tfrac12(1+P)q=q_0+q_2+q_4+\cdots,\notag\\P_{(1)}q&=\tfrac12(1-P)q=q_1+q_3+q_5+\cdots,\tag{16}\\Qq&=q_0+q_1-q_2-q_3+q_4+q_5-\cdots,\tag{17}\\Kq&=q_0-q_1-q_2+q_3+q_4-q_5-\cdots.\tag{18}\end{align}
Clearly, we might introduce $Q_{(0)}$, $Q_{(1)}$, and $K_{(0)}$, $K_{(1)}$, corresponding to $P_{(0)}$ and $P_{(1)}$. All the linities recently considered ($V_c$, $\phi$, $K\phi$) may be referred to as the capital linities.

\noindent\textsc{\S 2. Theorem of Independence and Allied Algebraic Features.}---If $q_1q_2=-q_2q_1$, then $q_1$ and $q_2$ are said to be anti-commutative. In the following enunciation, a product of $a$ different multenions from among $q_1$, $q_2$, \ldots{} $q_m$ is said to be odd-formed or even-formed, according as $a$ is odd or even.

\indent \emph{If $q_1$, $q_2$, \ldots{} $q_m$ be $m$ multenions such that $q_1^2$, $q_2^2$, \ldots{} $q_m^2$ are all scalars differing from zero, and each pair is anti-commutative; then, if $m$ is even, the $2^m$ multiplicative combinations are linearly independent; and, if $m$ is odd, they are independent unless the product $q_1q_2\ldots q_m$ of all of them is a scalar, and, if it is a scalar, $2^{m-1}$ and only $2^{m-1}$ of the combinations are independent, and these $2^{m-1}$ independent combinations may be taken as the odd-formed combinations, or instead as the even-formed combinations, or instead as the combinations from which one assigned multenion, such as $q_m$, is absent.} The theorem is true of linities as well as of multenions.

\indent Thus, if $n$ is even, a general multenion cannot be expressed by fewer than $2^n$ independent scalars. In the case of $n=4$, it cannot be expressed by fewer than sixteen independent scalars. In the case of $n=3$, it can be expressed by four independent scalars, the same number as in the case of $n=2$. [But to express it thus, we have to impose the condition that $i_1i_2i_3$ is a scalar such as $-1$. This is practically Hamilton's procedure in Quaternions.] Either of these cases might be regarded as the case of

\end{document}
